% =============================================================================
% preamble/glossary.tex: abbreviations and glossary (glossaries-extra).
%
% WHAT YOU GET
% Appendix A has two numbered sections, printed by
% backmatter/appendix-glossary.tex:
%   A.1 Abbreviations  short form -> long form   (e.g. MCP  Model Context Protocol)
%   A.2 Glossary       term -> definition
% Entries are defined in glossary/entries.tex. Only entries used somewhere
% in the text are printed, alphabetically. Sorting is done by LaTeX itself
% (\makenoidxglossaries), so no extra program is needed on Overleaf.
%
% USING ENTRIES IN THE TEXT (every use is a link to the appendix entry)
%   \gls{mcp}          first use "Model Context Protocol (MCP)", later "MCP"
%   \glspl{llm}        plural: first use "large language models (LLMs)"
%   \Gls{...} \Glspl   capitalised first letter, for the start of a sentence
%   \glsxtrshort{mcp}  always "MCP"
%   \glsxtrlong{mcp}   always "Model Context Protocol"
%   \glsxtrfull{mcp}   always "Model Context Protocol (MCP)"
%   \gls{def-mcp}      "Model Context Protocol", linking to its definition
%   \gls{rug-pull}     a glossary term, linking to its definition
%   \glsfmtshort{mcp}  "MCP" in a chapter/section title or caption; there, \gls
%                      would use up the first-use expansion and put links in
%                      the table of contents
%   \glsreset{mcp}     make the next \gls{mcp} expand again (e.g. per chapter)
%
% ABBREVIATION -> DEFINITION LINK
% For an abbreviation that also needs a definition, \newtermabbr (defined
% below) creates two linked entries. Its line in A.1 ends with
% "(see definition)", which links to the definition in A.2.
% =============================================================================

% --- Package options ---------------------------------------------------------------------
\usepackage[
  abbreviations,               % create the separate "abbreviations" list (A.1)
  nonumberlist,                % don't list page numbers after each entry
  section    = section,        % print each list as a \section, so both sit
                               % inside the appendix chapter
  numberedsection = autolabel, % number those sections (A.1, A.2) and give each
                               % a \label: \glsautoprefix + list name
  toc        = false,          % numbered sections reach the table of contents
                               % anyway; this avoids a duplicate entry
  nogroupskip,                 % no extra gap between letter groups (A, B, ...)
  stylemods,                   % load glossaries-extra's fixes to the list styles
                               % (needed for the "(see definition)" hook below)
]{glossaries-extra}

\renewcommand*{\glsautoprefix}{app:}  % section labels become app:abbreviations
                                      % (A.1) and app:main (A.2), for \cref
\makenoidxglossaries                  % sort and print lists without makeindex

% --- Appearance ----------------------------------------------------------------------------
% long-short: first use prints "long form (SHORT)", later uses "SHORT".
% Other styles: short-long ("SHORT (long form)"), long-short-sc (small caps).
\setabbreviationstyle{long-short}             % ordinary abbreviations
\setabbreviationstyle[termabbr]{long-short}   % abbreviations made by \newtermabbr
                                              % (their own category, see below)
\renewcommand*{\glsnamefont}[1]{\textbf{#1}}  % entry names in bold in both lists
\setlength{\glsdescwidth}{0.75\linewidth}     % width of the long-form column in
                                              % A.1 (the "long" list style)

% --- \newtermabbr: abbreviation + definition, linked -----------------------------------------
% Usage: \newtermabbr{label}{SHORT}{long form}{definition}
% Creates:
%   label      an abbreviation (A.1). Its "category" is set to termabbr so the
%              two hooks below apply only to these entries.
%   def-label  a glossary term (A.2) named "long form (SHORT)" with the
%              definition; in the text \gls{def-label} prints just "long form".
% #1..#4 are the four arguments in order.
\newcommand{\newtermabbr}[4]{%
  \newglossaryentry{def-#1}{name={#3 (#2)}, text={#3}, description={#4}}%
  \newabbreviation[category=termabbr]{#1}{#2}{#3}%
}

% Helper commands for the hooks: add the definition entry to the Glossary,
% and print the "(see definition)" link to it.
\newcommand*{\addtermdefinition}[1]{\glsadd{def-#1}}
\newcommand*{\linktermdefinition}[1]{%
  \space(\emph{see} \glshyperlink[definition]{def-#1})}

% Hook 1, runs after every \gls-type use of a termabbr abbreviation: also
% mark its definition as used, so the definition always appears in A.2 and
% the "(see definition)" link always has a target.
% \glslabel holds the label just used; \expandafter turns it into its text
% (e.g. "mcp") before the helper builds "def-mcp".
\glsdefpostlink{termabbr}{\expandafter\addtermdefinition\expandafter{\glslabel}}

% Hook 2, runs after the description of each termabbr entry in the printed
% list: append "(see definition)". \glscurrententrylabel is the label of the
% entry being printed.
\glsdefpostdesc{termabbr}{%
  \expandafter\linktermdefinition\expandafter{\glscurrententrylabel}}

% --- Entries -------------------------------------------------------------------------------
% Must come after the settings above.
% =============================================================================
% glossary/entries.tex: abbreviation and glossary entries.
% Loaded at the end of preamble/glossary.tex (see that file for all the commands you can use in the text).
%
% Only entries used in the text appear in Appendix A, sorted alphabetically, so defining an entry you never use is harmless. Keep this file grouped by kind and alphabetical by label so entries are easy to find.
%
% THE THREE KINDS OF ENTRY
%
% 1. Abbreviation only -> A.1 Abbreviations
%      \newabbreviation{label}{SHORT}{long form}
%    Write the long form in lower case unless it is a proper noun; \Gls capitalises it at the start of a sentence.
%      \newabbreviation[longplural={...}, shortplural={...}]{label}{SHORT}{long form}
%
% 2. Defined term only -> A.2 Glossary
%      \newglossaryentry{label}{name={term}, description={definition}}
%    Optional keys: plural={...} for an irregular plural; sort={...} to sort under a different spelling.
%
% 3. Abbreviation that also has a definition -> both lists, linked
%      \newtermabbr{label}{SHORT}{long form}{definition}
%    In the text, \gls{label} gives the abbreviation and \gls{def-label} the defined term. The A.1 line ends "(see definition)" and links to A.2.
%
% Labels: lower case, no spaces; hyphens are fine (e.g. rug-pull). The label is only used in \gls{...}, never printed. Commas or "=" inside a value need braces: description={a, b, and c}.
% =============================================================================

% --- 1. Abbreviations ----------------------------------------------------------------------


% --- 3. Abbreviations with a definition ------------------------------------------------------
% The fourth argument (the definition) may start on the next line, but there must be no blank line between it and the third argument.


% --- 2. Defined terms ------------------------------------------------------------------------


